\documentclass[11pt,a4paper]{article}
\usepackage[margin=23mm]{geometry}
\usepackage{fontspec}
\setmainfont{Noto Serif CJK JP}
\setsansfont{Noto Sans CJK JP}
\setmonofont{DejaVu Sans Mono}[Scale=0.80]
\XeTeXlinebreaklocale "ja"
\XeTeXlinebreakskip=0pt plus 0.05em minus 0.01em
\usepackage{amsmath,amssymb,mathtools}
\usepackage{booktabs,array}
\usepackage[unicode,hidelinks]{hyperref}
\urlstyle{tt}
\hypersetup{pdftitle={DBI 流れの二分枝連続族と Maxwell 極限},pdfauthor={DBI algebraic research note}}
\setlength{\parindent}{1em}
\setlength{\parskip}{3pt}
\setlength{\emergencystretch}{3em}
\linespread{1.04}
\setlength{\abovedisplayskip}{8pt plus 2pt minus 2pt}
\setlength{\belowdisplayskip}{8pt plus 2pt minus 2pt}
\newcommand{\sym}{\operatorname{sym}}
\newcommand{\diag}{\operatorname{diag}}
\newcommand{\Id}{\mathbf 1}
\newcommand{\Tmix}{\mathsf T}
\begin{document}
\begin{center}
{\Large\bfseries DBI 流れの二分枝連続族と Maxwell 極限}\par
\vspace{4pt}
{\large 同一混合応力を持つ実配置と結論の適用範囲}\par
\vspace{4pt}
2026 年 11 月 1 日
\end{center}
\noindent\textbf{要旨}\quad
四次元単一実スカラー DBI について、同じ完全な混合応力と同じラグランジアンを持つが、
固定場での結合定数微分 $Q$ が異なる二つの実配置を連続族として与える。
時間的勾配と空間的勾配を組み合わせたこの族は、共通の Maxwell 点へ任意に近づく。
したがって、両セクターを含むその点の全場・結合定数近傍に、単一の応力のみの自律関数
$Q=f(\Tmix)$ は存在しない。$\lambda$ を引数に残す流れ、固定した勾配符号の全領域、
一般の単純スペクトルに対する不可能性は示さない。新規性は未確認である。

\section{規約と固定場での微分}
$g=\diag(-1,1,1,1)$ とし、$v=(\partial_\mu\phi)$ は共変勾配の列、
$F^{\mathsf T}=-F$ とする。正の平方根を取り、$\sym A=(A+A^{\mathsf T})/2$ と置く。
\begin{equation}\label{eq:def}
\begin{gathered}
h=g+\lambda vv^{\mathsf T},\qquad G=h+\sqrt\lambda F,\qquad
\Delta=\sqrt{\det G/\det g}>0,\\
L=\frac{1-\Delta}{\lambda},\qquad
Q=\left.\frac{\partial L}{\partial\lambda}\right|_{g,F,v},\qquad
X=v^{\mathsf T}g^{-1}v .
\end{gathered}
\end{equation}
応力は元の DBI 作用の Hilbert 応力を用い、改良項を加えない。
反変成分の行列 $T$ と混合成分の行列 $\Tmix$ を区別する。
\begin{equation}\label{eq:stress}
T=\frac{g^{-1}-\Delta\sym(G^{-1})}{\lambda},\qquad
\Tmix=Tg,\qquad U=\Id-\lambda\Tmix .
\end{equation}
比較には共通の固定質量尺度 $\mu$ を用い、
\begin{equation}\label{eq:units}
\widehat\lambda=\lambda\mu^4,\quad
(\widehat v,\widehat F)=(v,F)/\mu^2,\quad
(\widehat T,\widehat L)=(T,L)/\mu^4,\quad
\widehat Q=Q/\mu^8
\end{equation}
と無次元化した後、帽子を省略する。状態ごとに単位や結合定数の規格化を変えない。

\section{二つの制限された軸方向配置}
独立な電磁場成分は $F_{01}=e>0$ のみとする。
時間的配置 A と空間的配置 B を次のように定める。
\begin{equation}\label{eq:axial}
\begin{array}{c|c|c|c}
 &v& r=\det h/\det g&P=\Delta^2\\ \hline
\mathrm A &(a,0,0,0)^{\mathsf T}&1-\lambda a^2&1-\lambda(a^2+e^2)\\
\mathrm B &(0,b,0,0)^{\mathsf T}&1+\lambda b^2&1+\lambda(b^2-e^2)
\end{array}
\end{equation}
この ansatz では、行列の逆を取るだけで
\begin{equation}\label{eq:U}
U_{\mathrm A}=\diag(\Delta^{-1},r\Delta^{-1},\Delta,\Delta),\qquad
U_{\mathrm B}=\diag(r\Delta^{-1},\Delta^{-1},\Delta,\Delta)
\end{equation}
を得る。以後の構成はこの二つの制限された配置内で行う。

\newpage
\section{同一応力を持つ連続族}
任意の
\begin{equation}\label{eq:domain}
q>0,\qquad \rho>s>q
\end{equation}
を選ぶ。各行では場 $v,F$ を固定したまま $\lambda$ を変え、指定した端点
$\lambda_{\mathrm A}$ または $\lambda_{\mathrm B}$ で物理量を評価する。
\begin{center}
\renewcommand{\arraystretch}{1.65}
\begin{tabular}{@{}lcc@{}}
\toprule
 & 時間的配置 A & 空間的配置 B\\
\midrule
端点の結合定数 & $\displaystyle\lambda_{\mathrm A}=\frac{\rho-q}{\rho q}$
 & $\displaystyle\lambda_{\mathrm B}=\frac{s-q}{sq}$\\
勾配の振幅 & $\displaystyle a^2=\frac{q(\rho-s)}{\rho}$
 & $\displaystyle b^2=\frac{q(\rho-s)}{s}$\\
電場の振幅 & $\displaystyle e_{\mathrm A}^{\,2}=\frac{q(q+s)}{\rho}$
 & $\displaystyle e_{\mathrm B}^{\,2}=\frac{q(\rho+q)}{s}$\\
勾配の符号 & $X_{\mathrm A}=-a^2<0$ & $X_{\mathrm B}=b^2>0$\\
$\Delta$ & $q/\rho$ & $q/s$\\
$L$ & $q$ & $q$\\
$Q$ & $\rho q/2$ & $sq/2$\\
\bottomrule
\end{tabular}
\end{center}
$a,b,e_{\mathrm A},e_{\mathrm B}$ は表示した正数の正の平方根とする。
式 \eqref{eq:stress} から得られる完全な混合応力は
\begin{equation}\label{eq:same-T}
\boxed{\Tmix_{\mathrm A}=\Tmix_{\mathrm B}=\diag(-\rho,-s,q,q),\qquad
Q_{\mathrm A}-Q_{\mathrm B}=\frac{q(\rho-s)}2>0 .}
\end{equation}
従って、応力の全不変量を使っても二つの $Q$ を区別できず、$L$ を加えても区別できない。
これは結合定数が異なる二点の比較であり、同じ $\lambda$ での非一意性ではない。

$Q$ の計算では、表の端点の式を先に $L$ に代入して微分してはならない。
$P(\lambda)$ を式 \eqref{eq:axial} の固定場での一次関数として、
\begin{equation}\label{eq:Q-direct}
Q(\lambda)=\frac{\Delta-1}{\lambda^2}
-\frac{P'(\lambda)}{2\lambda\Delta}
\end{equation}
を求めてから、それぞれの端点を代入する。
この直接微分と行列式・逆行列の評価が、表と式 \eqref{eq:same-T} の検算である。

\section{弱結合から端点まで実枝に留まること}
各配置について $0\leq\lambda\leq\lambda_{\mathrm A,\mathrm B}$ を考える。
$P$ は一次関数であり、その両端は
\begin{equation}\label{eq:positive}
P_{\mathrm A}(0)=P_{\mathrm B}(0)=1,\qquad
P_{\mathrm A}(\lambda_{\mathrm A})=(q/\rho)^2>0,\qquad
P_{\mathrm B}(\lambda_{\mathrm B})=(q/s)^2>0 .
\end{equation}
したがって区間全体で $P>0$ である。また
$r_{\mathrm A}=P_{\mathrm A}+\lambda e_{\mathrm A}^2>0$、
$r_{\mathrm B}=1+\lambda b^2\geq1$ である。
よって両配置で $h$ は Lorentz 計量、$G$ は非退化であり、正の平方根の実 DBI 枝に沿って
$\lambda=0$ から端点まで連結する。これは代数方程式の候補根だけを並べた反例ではない。

\newpage
\section{共通の Maxwell 点へ近づく族}
$q>0$ を固定し、
\begin{equation}\label{eq:eps}
\rho=q+2\varepsilon,\qquad s=q+\varepsilon,\qquad\varepsilon>0
\end{equation}
と置く。表の二配置は、それぞれ
\begin{equation}\label{eq:weak}
\begin{gathered}
\lambda_{\mathrm A}\longrightarrow0^+,\qquad
\lambda_{\mathrm B}\longrightarrow0^+,\qquad
v_{\mathrm A},v_{\mathrm B}\longrightarrow0,\qquad
 e_{\mathrm A}^2,e_{\mathrm B}^2\longrightarrow2q,\\
\Tmix_{\mathrm A}=\Tmix_{\mathrm B}\longrightarrow\diag(-q,-q,q,q),\qquad
L_{\mathrm A}=L_{\mathrm B}=q,\qquad
Q_{\mathrm A},Q_{\mathrm B}\longrightarrow q^2/2
\end{gathered}
\end{equation}
を満たす。一方、任意の $\varepsilon>0$ では
\begin{equation}\label{eq:weak-gap}
\Tmix_{\mathrm A}=\Tmix_{\mathrm B}=\diag(-q-2\varepsilon,-q-\varepsilon,q,q),
\qquad Q_{\mathrm A}-Q_{\mathrm B}=q\varepsilon/2>0 .
\end{equation}
両者は同じ $F_{01}=\sqrt{2q}$、$v=0$、$\lambda=0$ の Maxwell 点へ近づく。
$q>0$ を固定するため非零電磁場への弱結合極限であり、弱電場や真空への極限ではない。
両方の $Q$ とその差は連続的に収束し、不連続性を主張しているのでもない。

\paragraph{近傍に対する結論}
場と結合定数を座標とする空間で、この Maxwell 点の任意の全近傍を取る。
 $\lambda\geq0$ の相対位相での近傍を $\lambda>0$ に制限してもよい。式 \eqref{eq:weak} より、十分小さい
$\varepsilon>0$ で A と B の両点がその近傍に入る。
もし近傍全体に単一の関数 $Q=f(\Tmix)$ が存在すれば、同じ入力
$\Tmix_{\mathrm A}=\Tmix_{\mathrm B}$ に対して同じ値を返す必要があるが、
式 \eqref{eq:weak-gap} と矛盾する。従ってこのような関数は存在しない。
この議論には $f$ の連続性や解析性さえ仮定しない。

ただし、時間的・空間的の\textbf{両セクターを含む全近傍}という量化が重要である。
一方の符号だけに制限した集合や、片方の族を除いた部分領域への結論ではない。
また極限点そのものでは二つの $Q$ は一致する。不一致は、極限へ任意に近い二点で生じる。

\paragraph{二種類の経路を区別する}
前節の実枝の確認では、各配置の場を固定して $\lambda$ だけを変えた。
これに対して式 \eqref{eq:eps} の極限では、$\lambda$ とともに $v,F$ も変わる。
ここで比較した $Q$ は、各 $\varepsilon$ の配置における
$\partial_\lambda L|_{g,F,v}$ であり、この族に沿う全微分 $dL/d\varepsilon$ ではない。
$L=q$ が族に沿って一定であることは、$Q=0$ を意味しない。

\section{各 ansatz 内での分枝選択}
式 \eqref{eq:U} は、A では $U_{00}U_{22}=1$、B では $U_{11}U_{22}=1$ を与える。
目標応力 \eqref{eq:same-T} に対して、対応する条件は
\begin{equation}\label{eq:pair}
\begin{aligned}
\mathrm A:&\quad(1+\lambda\rho)(1-\lambda q)=1
 &&\Longrightarrow\quad\lambda=\frac{\rho-q}{\rho q},\\
\mathrm B:&\quad(1+\lambda s)(1-\lambda q)=1
 &&\Longrightarrow\quad\lambda=\frac{s-q}{sq}
\end{aligned}\qquad(\lambda\ne0)
\end{equation}
であり、\textbf{各制限 ansatz 内}では非零の $\lambda$ が一意に選ばれる。
異なる配置で逆数対の割り当てが変わることが、同一応力で二つの値が生じる仕組みである。
これは一般の固定 $X$ 符号セクターにおける一意性定理ではない。

\newpage
\section{既存結果との関係}
本ノートは DBI 行列式、Hilbert 応力、逆数対を持つ特性因数分解そのものを新結果としない。
固定原稿の \texttt{eq:char-factor-app}、\texttt{eq:beta-app} と、前段の二点反例を出発点としている。
今回の追加は、反例を連続族にし、共通の弱結合点へ近づけたうえで、結論の範囲を明示したことである。

\paragraph{四次元 DBI の流れに対する問い}
Brizio らの \href{https://arxiv.org/abs/2602.24058v2}{arXiv:2602.24058v2} 第 5.2 節の式 (119) の後は、
四次元単一スカラー DBI の応力テンソル流れを $f(T;\lambda)$ の形で問うている。
この問いは、$\lambda$ を陽に含まない自律形 $f(T)$ だけを要求するものではない。
本族では $\lambda_{\mathrm A}\ne\lambda_{\mathrm B}$ なので、応力だけの自律化への反例は
$\lambda$ を残す流れの存在に対する反例ではない。

\paragraph{純粋な非線形電磁気学との区別}
\href{https://arxiv.org/abs/2610.35135v1}{arXiv:2610.35135v1} 第 4 節の局所・大域の議論は、
純粋な非線形電磁気学の文脈にある。スカラー勾配が非零の二つの配置を比較する本構成へ、
そのまま定理を移してはいない。
また \href{https://arxiv.org/abs/2309.04253v4}{arXiv:2309.04253v4} 第 5 節の定理 3 は
$\lambda$ 依存を許す枠組みであり、
\href{https://arxiv.org/abs/2610.33887v1}{arXiv:2610.33887v1} 第 2.2--2.3 節の分枝可逆性の議論も、
どの枝と領域を固定するかを区別して比較する必要がある。
これらは範囲の比較であって、本連続族の新規性を確定する文献調査ではない。

\section{検算記録と到達点}
同梱の \texttt{check.py} は、$s=q+u$、$\rho=q+u+w$、$q,u,w>0$ と置いて
仮定 \eqref{eq:domain} を実装している。行列式と逆行列を直接評価し、
\texttt{results.json} に\textbf{成分ごとの厳密零残差 50 件}を記録した。
内訳は、各配置の $\Delta^2$、混合応力 16 成分、$L$、$Q$ の計 38 件と、
逆数対 2 件、流れの差 1 件、弱結合極限 8 件、弱結合族の流れの差 1 件である。
区間正値性はこれらとは別に、一次性と正の端点に対する assertion で検査した。

記録環境は Python 3.12.14、SymPy 1.14.0、CPU 実行時間 0.363 秒であり、
探索・数値近似を必要としない。一回の環境依存の時間であって性能保証ではない。
スクリプトは文献の内容や新規性を検証するものではない。
別実装の \texttt{independent\_check.py} でも、端点の代入前に対角計量成分の直接変分と
固定場の結合定数微分を行い、成分ごとの厳密零残差 43 件を確認した。
二つの検算には重複があり、件数は独立な物理恒等式の数ではない。

\paragraph{ここまでで示したこと}
共通の結合定数規格化と元の Hilbert 応力のもとで、両勾配符号を含む実 DBI 領域には
単一の $Q=f(\Tmix)$ が存在しない。さらに、その障害は Maxwell 点の任意の全近傍に残る。
ただし固有値 $q$ は二重であり、一般の単純スペクトルでの不可能性は示していない。
分枝局所的な自律化、固定 $X$ 符号の一般配置、改良応力や再パラメータ化は別問題である。
局所的一定場の代数的構成に留まり、動力学的安定性、因果的安定性、有限全エネルギーは主張しない。

\vspace{3pt}
\noindent\textbf{固定出典}\quad
\href{https://github.com/ryos-physrockme/dbi-stress-tensor-flow/tree/45c11bc3f2396bcb523a6e4906bba89316e16ce9}{\texttt{ryos-physrockme/dbi-stress-tensor-flow}}\par
\noindent\small
取得基点コミット\quad\texttt{45c11bc3f2396bcb523a6e4906bba89316e16ce9}\par
\noindent
\href{https://github.com/ryos-physrockme/dbi-stress-tensor-flow/blob/45c11bc3f2396bcb523a6e4906bba89316e16ce9/paper/paper.tex}{参照原稿}と
\href{https://github.com/ryos-physrockme/dbi-stress-tensor-flow/blob/45c11bc3f2396bcb523a6e4906bba89316e16ce9/research/ch_autonomy_20261101/note.tex}{前段ノート}をこのコミットに固定した。
本ノートは追加検算の記録であり、新規性の確定や一般分類の完成を主張しない。
\end{document}
